\subsection{麦克莱恩可分性判别法}

定理 1. --- 设 K 是特征为 0 的域。每一个既约 K-代数，特别地 K 的每一个扩张，都在 K 上可分。

我们首先证明 K 的每个扩张 L 都在 K 上可分。设 B 是 L 在 K 上的一个超越基（V，第 109 页，定理 1），并令 \(L_{1} = \mathbf{K}(\mathbf{B})\) 。则 \(L_{1}\) 在 K 上可分（V，第 121 页，命题 6）。进一步，L 是 \(L_{1}\) 的代数扩张，且 \(L_{1}\) 是特征为 0 的域；因此 L 在 \(L_{1}\) 上可分（V，第 37 页，推论）。于是由命题 9，L 在 K 上可分。

现在设 A 是域 K 上的一个既约代数。根据命题 2 的推论（V，第 119 页），存在一个由 K 的扩张组成的族 \((L_{i})_{i \in I}\) ，使得 A 同构于 \(\prod_{i \in I} L_{i}\) 的一个子代数。由前所述，每个代数 \(L_{i}\) 在 K 上都是可分的，

因此由命题 3 a) 和 c)（V，第 119 页）可知 A 也具有同样的性质。

定理 2. --- 设 K 是特征为 \(p \neq 0\) 的域， \(\mathbf{K}^{p^{-\infty}}\) 是 K 的一个完全闭包，A 是一个交换 K-代数。以下性质等价：

\begin{enumerate}
\def\labelenumi{\alph{enumi})}
\item
  A 是可分的。
\item
  存在一个扩张 \(\mathbf{K}'\) （\(\mathbf{K}\) 的扩张），使得 \(\mathbf{K}'\) 是完全域，并且 \(\mathbf{K}' \otimes_{\mathbf{K}} \mathbf{A}\) 是既约的。
\item
  环 \(\mathbf{K}^{p^{-\infty}}\otimes_{\mathbf{K}}\mathbf{A}\) 是既约的。
\item
  环 \(\mathbf{K}' \otimes_{\mathbf{K}} \mathbf{A}\) 对每个扩张 \(\mathbf{K}'\) （\(\mathbf{K}\) 的扩张）都是既约的，只要该扩张是有限次的 \(p\) -根式扩张且高度 \(\leqslant 1\) 。
\item
  对于 A 中在 K 上线性自由的每一个元素族 \((a_{i})_{i\in I}\) ，族 \((a_{i}^{p})_{i\in I}\) 在 K 上是线性自由的。
\item
  存在向量 K-空间 A 的一组基 \((a_{i})_{i\in I}\) ，使得族 \((a_{i}^{p})_{i\in I}\) 在 K 上是线性自由的。
\end{enumerate}

如果扩张 \(\mathbf{K}'\) （\(\mathbf{K}\) 的扩张）是一个完全域，它就包含一个子扩张，该子扩张 \(\mathbf{K}\) -同构于 \(\mathbf{K}^{p^{-\infty}}\) （V，第 6 页，命题 3）；此外，每个 \(p\) -根式扩张（\(\mathbf{K}\) 的扩张）都同构于 \(\mathbf{K}^{p^{-\infty}}\) 的一个子扩张（V，第 26 页，命题 3）。这些评注表明了蕴含关系 \(a) \Rightarrow b) \Rightarrow c) \Rightarrow d)\) 。

我们来证明 \(d)\) 蕴含 \(e)\) 。设 \((a_i)_{i \in I}\) 是 \(A\) 中的一个线性自由族，并令 \((\lambda_i)_{i \in I}\) 是 \(K\) 中的一个有限支撑族，使得 \(\sum \lambda_i a_i^p = 0\) 。设 \(K'\) 为 \(K^{p^{-\infty}}\) 的由元素 \(\lambda_i^{p^{-1}}\) 生成的子扩张；它是有限次的且高度 \(\leqslant 1\) 。令 \(x = \sum_{i \in I} \lambda_i^{p^{-1}} \otimes a_i\) 在 \(K' \otimes_K A\) 中；我们有

\[
x ^ {p} = \sum_ {i \in I} \lambda_ {i} \otimes a _ {i} ^ {p} = 1 \otimes \sum_ {i \in I} \lambda_ {i} a _ {i} ^ {p} = 0.
\]

根据假设 \(d)\) 我们有 \(x = 0\) ，从而 \(\lambda_i = 0\) 对所有 \(i \in I\) 。

显然 \(e)\) 蕴含 \(f)\) ，并且只需证明 \(f)\) 蕴含 \(a)\) 。于是设 \((a_i)_{i \in I}\) 是 \(A\) 在 \(K\) 上的一组基，使得族 \((a_i^p)_{i \in I}\) 在 \(K\) 上线性自由。设 \(L\) 是 \(K\) 的一个扩张，并且令 \(x\) 为 \(L \otimes_K A\) 中满足 \(x^p = 0\) 的一个元素。记 \(x = \sum_{i \in I} \lambda_i \otimes a_i\) ，其中 \(\lambda_i \in L\) 对所有 \(i \in I\) 。我们有 \(x^p =\)

\(\sum_{i\in I}\lambda_i^p\otimes a_i^p = 0\) ，并且由于族 \((a_i^p)_{i\in I}\) 在 \(\mathbf{K}\) 上线性自由，我们有

\(\lambda_{i}^{p} = 0\) ，从而 \(\lambda_{i} = 0\) 对所有 \(i \in I\) ；由此可得 \(x = 0\) 。于是我们证明了 \(x^{p} = 0\) 蕴含 \(x = 0\) （在 \(\mathbf{L} \otimes_{\mathbf{K}} \mathbf{A}\) 中），由此立即得出 \(\mathbf{L} \otimes_{\mathbf{K}} \mathbf{A}\) 是既约的。

推论 1（麦克莱恩）。--- 设 K 是特征指数为 p 的域，Ω 是 K 的一个完全扩张，L 是 Ω 的一个子扩张。则以下条件等价：

\begin{enumerate}
\def\labelenumi{\alph{enumi})}
\item
  L在K上可分。
\item
  L 与 \(\mathbf{K}^{p^{-\infty}}\) 在 \(\mathbf{K}\) 上线性不相交。
\item
  L 在 K 上与包含于 \(\Omega\) 中的每个有限次且高度 \(\leqslant 1\) 的 p-根式扩张线性不相交。
\end{enumerate}

当 K 的特征为 0 时，情形是平凡的，因为此时 L 在 K 上可分（定理 1），且按约定 \(K^{p^{-\infty}} = K\) 。现假设 \(p \neq 1\) ，并首先证明 a) 蕴含 b)。假设 L 在 K 上可分，并设 \((a_i)_{i \in I}\) 为 L 在 K 上的一组基。设 \((\lambda_i)_{i \in I}\) 是 \(K^{p^{-\infty}}\) 中元素的一个有限支撑族，使得 \(\sum_{i \in I} \lambda_i a_i = 0\) ；则存在一个整数 \(f \geqslant 0\) 以及元素

\(\mu_{i}\) 属于 \(\mathbf{K}\) ，使得 \(\lambda_{i} = \mu_{i}^{p^{-f}}\) 。我们有

\[
\sum_ {i \in I} \mu_ {i} a _ {i} ^ {p ^ {f}} = \left(\sum_ {i \in I} \lambda_ {i} a _ {i}\right) ^ {p ^ {f}} = 0
\]

且定理 2 通过对 \(f\) 进行归纳推出，族 \((a_i^{p^f})_{i\in I}\) 在 \(K\) 上线性自由。因此我们有 \(\mu_i = 0\) ，从而 \(\lambda_i = 0\) 对所有 \(i\in I\) 。最后 \(L\) 与 \(K^{p^{-\infty}}\) 在 \(K\) 上线性不相交。

显然 \(b)\) 蕴含 \(c)\) 。最后假设 \(c)\) 成立，并令 \(K'\) 是 \(K\) 的一个扩张，它是有限次的 \(p\) -根式扩张且高度 \(\leqslant 1\) 。环 \(K' \otimes_K L\) 同构于 \(\Omega\) 的一个子环，因此是既约的。现在由定理 2 可知 \(L\) 在 \(K\) 上可分。

推论 2. --- 设 K 为特征指数为 p 的域， \(K^{p^{-\infty}}\) 为 K 的一个完全闭包，L 为 K 的一个可分扩张。则环 \(L \otimes_{K} K^{p^{-\infty}}\) 是一个域。此外，若 L 在 K 上是代数的，则 \(L \otimes_{K} K^{p^{-\infty}}\) 是 L 的一个完全闭包。

情形 \(p = 1\) 是平凡的，因此我们假设 \(p \neq 1\) 。设 \(\Omega\) 是 \(L\) 的一个完全闭包。根据推论 1，存在一个 K-代数同态，把 \(L \otimes_K K^{p^{-\infty}}\) 映到 \(L[K^{p^{-\infty}}]\) 上，它将 \(x \otimes y\) 映为 \(xy\) （对于 \(x \in L\) 和 \(y \in K^{p^{-\infty}}\) ）。由于 \(K^{p^{-\infty}}\) 在 \(K\) 上是代数的，环 \(L[K^{p^{-\infty}}]\) 是 \(\Omega\) 的一个子域（V，第 18 页，推论 1）。进一步假设 \(L\) 在 \(K\) 上是代数的，则 \(L[K^{p^{-\infty}}]\) 是完全域 \(K^{p^{-\infty}}\) 的代数扩张，因此它是一个完全域（V，第 43 页，推论 1）；最后，由于域 \(L[K^{p^{-\infty}}]\) 是一个 \(p\) -根式扩张（\(L\) 的扩张）（V，第 25 页，推论），它是 \(L\) 的一个完全闭包。

推论 3. --- 设 L 是 K 的一个代数扩张，M 是一个交换 L-代数（例如 L 的一个扩张）。若 M 在 K 上可分，则它在 L 上可分。

代数 L 与 M 的一个子代数 K-同构，因此 L 是 K 的一个可分扩张。因此（推论 2）存在一个 L-同构，将 \(L^{p^{-\infty}}\) 映到 \(K^{p^{-\infty}} \otimes_{K} L\) 上。环 \(L^{p^{-\infty}} \otimes_{L} M\) 同构于 \((\mathrm{K}^{p^{-\infty}} \otimes_{K} L) \otimes_{L} M\) ，从而同构于 \(K^{p^{-\infty}} \otimes_{K} M\) （II，第 278 页，命题2），且后一个环是既约的，因为 M 在 K 上可分。因此环 \(L^{p^{-\infty}} \otimes_{L} M\) 是既约的，这证明了 M 在 L 上可分（V，第 122 页，定理 2）。

注记. --- 麦克莱恩判别法的表述可以不引入除 L 以外的 K 的任何扩张。事实上，由推论 1 的 c) 可知，L 在 K 上可分当且仅当 L 与 \(K^{p^{-1}}\) 在 K 上线性不相交。由于映射 \(x \mapsto x^{p}\) 是 L 到子域 \(L^{p}\) 上的同构，我们得到以下判别准则（参见 V，第 177 页，习题 11，其中给出了代数情形的类似准则）：

L 在 K 上可分当且仅当 L 的子域 \(L^{p}\) 与 K 在 \(K^{p}\) 上线性不相交。

